\documentclass{article}
\usepackage[utf8]{inputenc}
\usepackage[german]{babel}
\usepackage{amsmath, amssymb, amsthm}
\usepackage{booktabs}
\usepackage{array}

\newcolumntype{L}[1]{>{\raggedright\arraybackslash}p{#1}}

\title{Die fundamentale Symmetrie zwischen $\forall$ und $\bigcap$}
\author{Nach Daniel J. Velleman: \emph{How to Prove It}}
\date{\today}

\begin{document}

\maketitle

\section*{Die Brücke zwischen Logik und Mengenlehre}

Deine Beobachtung ist absolut exakt. Schauen wir uns die Definition des Durchschnitts einer Familie an. Ein Element $x$ liegt genau dann im Durchschnitt, wenn es in \emph{jedem einzelnen} $U_i$ liegt:
\[
x \in \bigcap_{i \in I} U_i \iff \forall i \in I \, (x \in U_i)
\]
Hier sieht man sofort: Die Operation $\bigcap$ wird auf der rechten Seite direkt durch das logische Symbol $\forall$ definiert!

---

\section*{Gegenüberstellung im Scratch-Format}

Wenn wir Vellemans Rezepte anwenden, verhalten sich die Aussage über die Zugehörigkeit zur Topologie und die Aussage über den Durchschnitt eines Elements nahezu identisch:

\begin{center}
\small
\begin{tabular}{ L{0.45\textwidth} | L{0.45\textwidth} }
\toprule
\textbf{Die logische Eigenschaft} & \textbf{Das mengentheoretische Objekt} \\
\midrule
\textbf{Aussage:} \newline
$\forall i \in I \, (U_i \in \mathcal{T})$ \newline
\newline
(\emph{Bedeutung: Alle Mengen der Familie liegen einzeln in der Topologie.}) &
\textbf{Aussage:} \newline
$x \in \bigcap_{i \in I} U_i \iff \forall i \in I \, (x \in U_i)$ \newline
\newline
(\emph{Bedeutung: Ein konkretes Element $x$ liegt in allen Mengen der Familie gleichzeitig.}) \\
\hline
\textbf{Rezept-Anwendung (Given):} \newline
Wenn du ein konkretes $i_0 \in I$ hast, darfst du eintragen: \newline
$\rightarrow U_{i_0} \in \mathcal{T}$ &
\textbf{Rezept-Anwendung (Given):} \newline
Wenn du dasselbe konkrete $i_0 \in I$ hast, darfst du eintragen: \newline
$\rightarrow x \in U_{i_0}$ \\
\hline
\textbf{Rezept-Anwendung (Goal):} \newline
Um die Aussage zu beweisen, musst du sagen: \newline
$\rightarrow$ Sei $i_0 \in I$ beliebig. \newline
$\rightarrow$ \emph{Neues Goal:} $U_{i_0} \in \mathcal{T}$ &
\textbf{Rezept-Anwendung (Goal):} \newline
Um zu beweisen, dass $x$ im Durchschnitt liegt, musst du sagen: \newline
$\rightarrow$ Sei $i_0 \in I$ beliebig. \newline
$\rightarrow$ \emph{Neues Goal:} $x \in U_{i_0}$ \\
\bottomrule
\end{tabular}
\end{center}

\section*{Fazit der Analogie}

Es gilt die wunderbare mathematische Eselsbrücke:
\begin{itemize}
    \item Der \textbf{Durchschnitt} ($\bigcap$) ist nichts anderes als ein riesiges, eventuell unendliches \textbf{UND} ($\land$). Er sammelt Elemente, die in \emph{allen} Mengen gleichzeitig sein müssen.
    \item Die \textbf{Vereinigung} ($\bigcup$) ist das Gegenstück: Ein riesiges \textbf{ODER} ($\lor$). Sie sammelt Elemente, die in \emph{mindestens einer} Menge sein müssen.
\end{itemize}

\end{document}
